\begin{lema}[Reducción a torres en el último tablero]\label{terminal-pivotes}
Sea $A$ la máscara del tablero $B^{(K)}$. En $F\cap D_K^c$ hay una
colocación que completa el tablero y
\[
 Q(B^{(K)})=\per A,\qquad L(B^{(K)})=L_T(A),
\]
donde $L_T(A)=H_T(p)+\log\Gamma_2$ es el predictor de torres en el mismo soporte, como en el modelo
de torres con máscara de~\cite[Secciones~8.1 y~8.5]{pred}. Aquí $H_T(p)=-\sum_e p_e\log p_e
+\sum_e(1-p_e)\log(1-p_e)$ es la entropía de torres, evaluada en su punto estacionario con masas
de fila y de columna iguales a uno, y $\Gamma_2$ es el factor de la
Sección~\ref{costes:permanente}.
\end{lema}
\begin{proof}
En $F$ el punto de Bethe $p$ de $B^{(K)}$ tiene masas de fila y de columna iguales a uno y es positivo en
todas las casillas disponibles. Para un conjunto cualquiera $I$ de filas, toda su masa $|I|$ llega al
conjunto
de columnas vecinas, que admite a lo sumo una unidad por columna. En consecuencia $|\mathcal
N(I)|\ge|I|$. Para ver directamente que se puede cubrir toda fila, tómese un emparejamiento de
cardinal máximo. Si deja una fila sin cubrir, desde todas las filas descubiertas se recorren
alternativamente aristas fuera del emparejamiento desde las filas y aristas del emparejamiento desde
las columnas. Una columna descubierta alcanzable daría un camino aumentante, contradicción. Luego cada
columna alcanzable está emparejada con una fila alcanzable, y estas filas emparejadas no agotan las
filas alcanzables, entre las que hay una descubierta. Todas las vecinas de una fila alcanzable son
alcanzables, incluida, si existe, la columna por la que se entró en esa fila. Esto contradice la
desigualdad para el conjunto de filas alcanzables. Así $\per A>0$.

En $D_K^c$, cada diagonal contiene a lo sumo un cruce vacío. Toda permutación del
soporte respeta por ello las dos restricciones diagonales. Además, las casillas de diagonales ocupadas
ya se excluyeron al formar $A$, por lo que cualquier permutación del soporte da una completación legal
de $B^{(K)}$. La construcción es reversible y prueba $Q(B^{(K)})=\per A$.

En una diagonal libre con una casilla disponible $e$, la vacancia es $h_f=1-p_e=:q_e$. Cada casilla
tiene dos de estas diagonales y una diagonal libre sin casillas tiene $h_f=1$. Así en $H_B$ los
términos diagonales restan dos copias de $\sum_e q_e\log q_e$, y queda exactamente la entropía de
torres
\[
 H_T(p)=-\sum_e p_e\log p_e+\sum_e q_e\log q_e .
\]
La ecuación estacionaria del Lema~\ref{coste-vacancias} pasa a ser $p_eq_e=a_ib_j$ en toda casilla
disponible, con los mismos $a_i,b_j$. También se obtiene $H_B=H_T$
directamente como identidad de los funcionales en los puntos factibles del soporte: en toda diagonal
de una sola casilla la vacancia siempre es el complemento de la única entrada. De aquí que el punto
de Bethe sea también el punto estacionario del predictor de torres.

Queda verificar el normalizador, como en \cite[Sección~8.1]{pred}. Para una casilla de masa $p$
ponemos $q=1-p$, $\tau_d=pq/(1-dp)$ y $b_d=(1-dp)/q^d$. Al eliminar la coordenada de su primera
diagonal, el pivote en $S_4$ es $q+\tau_4$. Si $r$ indica las otras líneas de la casilla, la
contribución $\tau_4rr^T$ en el bloque restante se sustituye, por el complemento de Schur, por
$(\tau_4-\tau_4^2/(q+\tau_4))rr^T$. La segunda diagonal de esa casilla se elimina después con $d=3$.
En ambos pasos
\[
 q+\tau_d=\frac{q(1-(d-1)p)}{1-dp},\qquad
 \tau_d-\frac{\tau_d^2}{q+\tau_d}=\tau_{d-1},\qquad
 b_d(q+\tau_d)=b_{d-1}.
\]
Dado que la diagonal eliminada tiene una sola casilla, la operación afecta sólo a la contribución de
esa casilla. Una línea sin casillas tiene pivote uno y no afecta a las otras coordenadas. El modo cero
$(\mathbf1_R,-\mathbf1_C,0)$ tiene coordenada cero en cada línea eliminada. Por ello podemos calcular
el pseudodeterminante en una base ortonormal de su complemento que contenga los vectores de las
diagonales eliminadas. Allí la factorización ordinaria por los pivotes es exacta y el modo de filas y
columnas conserva su longitud $\sqrt{2K}$. El factor correspondiente del normalizador es
$\sqrt{2\pi}\sqrt{2K}=\sqrt{4\pi K}$ tanto antes como después. Como $e^{C_d}=\prod_e b_d(p_e)$, las
identidades prueban
$e^{C_4}\pdet S_4=e^{C_2}\pdet S_2$ tras completar las dos eliminaciones por casilla. En particular
$\Gamma_4=\Gamma_2$. Esto da $L(B^{(K)})=H_T(p)+\log\Gamma_2=L_T(A)$.
\end{proof}
